\begin{frame}{Capacidade de acionamento e fan-out}
    \begin{block}{Fan-out}
        É a quantidade máxima de entradas que uma saída consegue acionar sem
        violar os níveis de tensão ou os limites de corrente.
    \end{block}

    \begin{columns}[T]
        \begin{column}{0.48\textwidth}
            \begin{block}{Limite em corrente contínua}
                \[
                    \text{fan-out}_{H}=\frac{|I_{OH}|}{I_{IH}}
                \]
                \[
                    \text{fan-out}_{L}=\frac{I_{OL}}{|I_{IL}|}
                \]
            \end{block}
        \end{column}
        \begin{column}{0.48\textwidth}
            \begin{alertblock}{Limite dinâmico}
                Mesmo com baixa corrente estática, muitas entradas CMOS somam
                capacitância. Isso aumenta o tempo de subida, o tempo de descida
                e a corrente de comutação.
            \end{alertblock}
        \end{column}
    \end{columns}

    \centering
    O fan-out real é o menor limite obtido para os estados alto e baixo.
\end{frame}
